\documentclass{zhvt-classic}

% 使用本地源流明体（GenRyuMin2 TW）按文件路径加载，不依赖系统字体名
% 版框（内框／界栏／书口侧栏／鱼尾）改用与版心印章同一印泥棕色：
%   印章图片 yshs.png 实色 = RGB(135,68,52) = HTML 874434
\usepackage{xcolor}
\usepackage{tikz}
% sealbrown 双注册：xcolor 供 TikZ/pgf，L3 供 cls 版框
\definecolor{sealbrown}{HTML}{874434}
\ExplSyntaxOn
\color_set:nnn { sealbrown } { HTML } { 874434 }
\ExplSyntaxOff

\zhvtset{
  grid color       = sealbrown,
  main font        = GenRyuMin2TW-SB,
  main font path   = C:/Users/Administrator/texmf/fonts/opentype/public/GenRyuMin2TW/,
jiazhu font      = GenRyuMin2TW-M,
  jiazhu font path = C:/Users/Administrator/texmf/fonts/opentype/public/GenRyuMin2TW/,
  % 页码框容纳 4 字：正常页「九十七」+前置空 1 格=4 行恰好；
  % 100+ 页「一百零一」共 5 行超出，末字「一/二/三」被框裁掉（2026-10-04 实测），故放宽到 5
  shukou page number chars = 5,
  % 书口书名框高从默认 6 字减到 4 字，把省出的 2 字高拨给 chapter 框，
  % 使 6 字书名+卷名放进 9 字高的 chapter 框后不再压线（2026-10-05）
  shukou title chars = 4,
}

\title{史记}[殷本纪]
\author{汉·司马迁}
\maker{闲叟漫记·重製}


% ===== 索隐述赞样式（2026-11-08 四次定稿 · 愚大人亲手改定参数）=====
% ⚠ 竖排前提：自定义字体家族必须带 RawFeature={vertical:+vert}；
%   simhei.ttf 加 +vhal 会整体加载失败静默回退，只写 +vert；simfang 同理。
\newCJKfontfamily{\shuzanfont}[Path=C:/Windows/Fonts/,RawFeature={vertical:+vert}]{simhei.ttf}
\newCJKfontfamily{\shuzanbody}[Path=C:/Windows/Fonts/,RawFeature={vertical:+vert}]{simfang.ttf}
% 题头：152x30pt 圆角 6pt 棕底白字；baseline (0,2) 使整框落入界行格子不压版框线；
% 白字节点 x 不补偿、y +8.2pt（像素级墨迹实测 dx/dy≈0，L/R 留白约 20pt）。
\newcommand{\shuzanhead}[1]{%
  \begin{tikzpicture}[baseline={(0,2)}]
    \fill[sealbrown, rounded corners=6pt]
      (-76pt,-15pt) rectangle (76pt,15pt);
    \node[text=white, inner sep=0pt, anchor=center] at (0pt,8.2pt)
      {\shuzanfont\fontsize{23}{23}\selectfont #1};
  \end{tikzpicture}%
}
% ★ 修正书口卷名位置：chapter 框起始位置改回 cls 默认值，避免全局影响。
%   只在每个卷首页通过 \juan 命令里的空命令微调该页位置（安全，不会跨页污染）。
\ExplSyntaxOn
\AtBeginDocument{
  \dim_set:Nn \l__zhvt_shukou_chapter_title_start_dim
    { \l__zhvt_shukou_title_height_dim + \l__zhvt_font_size_dim * 2 }
  \dim_set:Nn \l__zhvt_shukou_page_number_start_dim
    { \l__zhvt_shukou_chapter_title_start_dim + \l__zhvt_shukou_chapter_title_height_dim }
}
\ExplSyntaxOff

\begin{document}

% ★★ 版心（书口）排布覆写：把书名放到鱼尾之下（2026-10-04）
%
%   cls 默认版心自上而下：书名（title 静态框，在鱼尾之上）→ 鱼尾 → 篇名（chapter 动态框）→ 页码
%   本模板改为：          鱼尾（位置不动）→ 书名 → 篇名 → 页码分割线 → 页码
%
%   手法：把 cls 放在 title 静态框里的书名，改放进 chapter 动态框（鱼尾之下）；
%         title 静态框留空。书名取 \l__zhvt_title_str，即 \title{...} 的第一参数，
%         所以换书时只改 \title，不必在此重复写书名。
%
%   ⚠ 三条必须遵守的写法（2026-10-04 实测）：
%     ① 汉字必须用 \hbox{...} 包住。flowfram 会把框内容装进一个宽 14.4mm（= 一个 line_sep
%        = 书口栏宽）的纵向 minipage 并按该宽度折行，而整页靠 PDF /Rotate 90 显示、
%        框内文字本身不变向 —— 不包 \hbox 时「围炉夜话」4 字会被逐字折成 4 行，
%        在版面上就成了横排 4 字并横向溢出版框。
%     ② 竖排下**一个 \hbox 就是一整列**。书名与篇名要同处一列，必须放进同一个 \hbox；
%        分成两个 \hbox 时，后一个会被当成新的一列、甩到隔壁栏去
%        （实测「序」跑到书口栏外侧 14.4mm 的正文列上）。中间的 \quad 在框内是横向空白，
%        显示出来正好是书名与篇名之间的纵向间隔。
%     ③ ★ 内容盒要写成 **与框等宽的定宽盒** `\hbox to \l__zhvt_shukou_chapter_title_height_dim{…\hfil}`。
%        实测：chapter 框是把内容盒沿纵向**居中**摆放的，内容宽度一变，书名就跟着漂：
%        自然宽度的 \hbox（书名 38.4mm）落在 x=157.2mm，加篇名后（57.6mm）落在 x=152.4mm，
%        目录页（无篇名）与正文页的书名就差半格。做成等宽盒后，无论有没有篇名，
%        内容固定从框顶起排，实测两页均落在 x=142.8mm，不再漂。
\ExplSyntaxOn
% 把书口 chapter 框（书名+篇名）起始位置下移 1 字高，避免 4 字卷名顶到鱼尾
% ⚠ cls 的 \__zhvt_calculate_dimensions: 在 \AtBeginDocument 里执行，会重算该 dim；
%   必须把覆写挂到 begin document 之后，否则会被二次覆盖（2026-10-05 实测）。
% (旧方案已移到导言区，直接修改 chapter 框起始位置)
% 卷标题命令 #1=书口篇名(鱼尾) #2=正文卷名（含卷序，卷名内空格用 ~）
%
% ★★ 目录页两项修订（2026-10-05，愚大人指出）★★
%   ① 页码起点统一对齐：
%      原状 —— 目录项写成 \par\hspace{4em}卷名\hspace{4em}页码，页码起点 = 卷名实宽 + 4em。
%      卷名实宽随字数变：「卷一~梁惠王上」6 字 +~、「卷七~离娄上」5 字 +~，
%      差 1 个字位 → 卷七~卷十 4 条的页码整体比其余 10 条高 27.21pt（一个字），看着上下不齐。
%      改法 —— 卷名补全角空格（宽 = 1 字 = 27.215pt）补足到 **7 字符位**再写进目录，
%      14 条页码起点遂全部落在同一条竖线上（起点对齐；末字随页码字数自然参差）。
%      ⚠ 只对目录里的副本 \l_tmpc_tl 补，正文卷名仍用 \l_tmpa_tl 原样 —— 正文像素不受影响。
%   ② 卷名整体上移半字位：前导 \hspace 由 4em 改 3.5em。
%      实测原状目录块在版框内偏下（天头留白 139.6pt / 地脚留白 110.7pt），
%      上移 0.5em = 13.6pt 后上下基本均分，视觉居中。
\NewDocumentCommand { \juan } { mm } {
	\clearpage
	\tl_set:Nn \l__zhvt_chapter_title_tl { #1 }
	\tl_set:Nn \l_tmpa_tl { #2 }
	\tl_set_eq:NN \l_tmpc_tl \l_tmpa_tl
	\int_step_inline:nn { 7 - \str_count:n { #2 } }
	  { \tl_put_right:Nn \l_tmpc_tl { 　 } }
	\tl_set:Nx \l_tmpb_tl { \zhnum { page } }
	\tl_put_left:Nn \l_tmpb_tl { \hspace { 4em } }
	\tl_put_left:Nn \l_tmpb_tl { \l_tmpc_tl }
	\tl_put_left:Nn \l_tmpb_tl { \par \hspace { 3.5em } }
	\addtocontents { toc } { \l_tmpb_tl }
	\skip_horizontal:n { \l__zhvt_font_size_dim * 4 }
	\l_tmpa_tl
	\par
	\thispagestyle { empty }
}
\cs_set_protected:Npn \__zhvt_set_shukou_content: {
	% ★ 鱼尾上方留白（title 静态框）嵌入闲叟钤印：
	%   整页 /Rotate 90，故此处把图用 \rotatebox{90} 转正为竖章；
	%   尺寸卡在版心条幅内（书内约 7mm 宽 × 33.8mm 高，留余量不压鱼尾）。
	\setstaticcontents * { title } {
		\vbox_to_ht:nn { \l__zhvt_baseline_skip_dim } {
			\vskip 14.2pt
			\hbox_to_wd:nn { \l__zhvt_shukou_title_height_dim } {
				\hfil
				\rotatebox [ origin = c ] { 90 } { \includegraphics [ width = 7mm , height = 33.8mm ] { yshs.png } }
				\hfil
			}
		}
	}
% ★ 版心（书口）统一改法（2026-11-16 愚大人定）：18pt＋框内居中＋书口文字同版框色
%   ① 原 27.2pt 时 7 字长篇名（如「廉颇蔺相如列传」）末字溢出 chapter 框、压住页码；
%      改 \fontsize{18pt}{18pt}，前后 \hfil 令字串在框内居中（不随篇名字数漂移）。
%   ② \vbox{\vskip 8.2pt} 定量下移，抵消缩字号后栏宽方向 8.2pt 上偏（像素实测 y 中心 543.3 = 原版）。
%      ⚠ \raisebox / \vspace* 在 flowfram 框内无效；换字号须同步换补偿：16pt→10.1、19pt→7.4。
%   ③ chapter 与 page~number 均置 \color{sealbrown}（HTML 874434），与版框／界行同色；正文仍黑。
%   ⚠ 页码字位、钤印、鱼尾、界行、正文一律未动；五部合订本同此口径。
	\setdynamiccontents * { chapter }
	{ \color{sealbrown} \vbox { \vskip 8.2pt \hbox to \l__zhvt_shukou_chapter_title_height_dim
		{ \hfil \fontsize{18pt}{18pt}\selectfont \l__zhvt_title_str \quad \l__zhvt_chapter_title_tl \hfil } } }
	\setdynamiccontents * { page~number } { \color{sealbrown} \quad \zhnum { page } }
}
\__zhvt_set_shukou_content:
\ExplSyntaxOff

\maketitle

% 单篇，无多卷目录，不渲染目录页




% 方案①：仅改 A 类错讹，头部与版式不变

\juan{殷本纪}{卷三~殷本纪第三}
\hspace*{2em}殷契。母曰简狄。有娀氏之女。为帝喾次妃。三人行浴。见玄鸟堕其卵。简狄取吞之。因孕生契。契长而佐禹治水有功。帝舜乃命契曰。﹁百姓不亲。五品不训。汝为司徒而敬敷五教。五教在宽。﹂封于商。赐姓子氏。契兴於唐。虞。大禹之际。功业著於百姓。百姓以平。\\
\hspace*{2em}契卒。子昭明立。昭明卒。子相土立。相土卒。子昌若立。昌若卒。子曹圉立。曹圉卒。子冥立。冥卒。子振立。振卒。子微立。微卒。子报丁立。报丁卒。子报乙立。报乙卒。子报丙立。报丙卒。子主壬立。主壬卒。子主癸立。主癸卒。子天乙立。是为成汤。\\
\hspace*{2em}成汤。自契至汤八迁。汤始居亳。从先王居。作帝诰。\\
\hspace*{2em}汤征诸侯。葛伯不祀。汤始伐之。汤曰。﹁予有言。人视水见形。视民知治不。﹂伊尹曰。﹁明哉。言能听。道乃进。君国子民。为善者皆在王官。勉哉。勉哉。﹂汤曰。﹁汝不能敬命。予大罚殛之。无有攸赦。﹂作汤征。\\
\hspace*{2em}伊尹名阿衡。阿衡欲奸汤而无由。乃为有莘氏媵臣。负鼎俎。以滋味说汤。致于王道。或曰。伊尹处士。汤使人聘迎之。五反然後肯往从汤。言素王及九主之事。汤举任以国政。伊尹去汤适夏。既丑有夏。复归于亳。入自北门。遇女鸠。女房。作女鸠女房。\\
\hspace*{2em}汤出。见野张网四面。祝曰。﹁自天下四方皆入吾网。﹂汤曰。﹁嘻。尽之矣。﹂乃去其三面。祝曰。﹁欲左。左。欲右。右。不用命。乃入吾网。﹂诸侯闻之。曰。﹁汤德至矣。及禽兽。﹂\\
\hspace*{2em}当是时。夏桀为虐政淫荒。而诸侯昆吾氏为乱。汤乃兴师率诸侯。伊尹从汤。汤自把钺以伐昆吾。遂伐桀。汤曰。﹁格女众庶。来。女悉听朕言。匪台小子敢行举乱。有夏多罪。予维闻女众言。夏氏有罪。予畏上帝。不敢不正。今夏多罪。天命殛之。今女有众。女曰。﹃我君不恤我众。舍我啬事而割政﹄。女其曰。﹃有罪。其奈何﹄。夏王率止众力。率夺夏国。众有率怠不和。曰。﹃是日何时丧。予与女皆亡﹄。夏德若兹。今朕必往。尔尚及予一人致天之罚。予其大理女。女毋不信。朕不食言。女不从誓言。予则帑僇女。无有攸赦。﹂以告令师。作汤誓。於是汤曰。﹁吾甚武﹂。号曰武王。\\
\hspace*{2em}桀败於有娀之虚。桀奔於鸣条。夏师败绩。汤遂伐三㚇。俘厥宝玉。义伯。仲伯作典宝。汤既胜夏。欲迁其社。不可。作夏社。伊尹报。於是诸侯毕服。汤乃践天子位。平定海内。\\
\hspace*{2em}汤归至于泰卷陶。中垒作诰。既绌夏命。还亳。作汤诰。﹁维三月。王自至於东郊。告诸侯群后。﹃毋不有功於民。勤力乃事。予乃大罚殛女。毋予怨。﹄曰。﹃古禹。皋陶久劳于外。其有功乎民。民乃有安。东为江。北为济。西为河。南为淮。四渎已修。万民乃有居。后稷降播。农殖百谷。三公咸有功于民。故後有立。昔蚩尤与其大夫作乱百姓。帝乃弗予。有状。先王言不可不勉。﹄曰。﹃不道。毋之在国。女毋我怨。﹄﹂以令诸侯。伊尹作咸有一德。咎单作明居。\\
\hspace*{2em}汤乃改正朔。易服色。上白。朝会以昼。\\
\hspace*{2em}汤崩。太子太丁未立而卒。於是乃立太丁之弟外丙。是为帝外丙。帝外丙即位三年。崩。立外丙之弟中壬。是为帝中壬。帝中壬即位四年。崩。伊尹乃立太丁之子太甲。太甲。成汤适长孙也。是为帝太甲。帝太甲元年。伊尹作伊训。作肆命。作徂后。\\
\hspace*{2em}帝太甲既立三年。不明。暴虐。不遵汤法。乱德。於是伊尹放之於桐宫。三年。伊尹摄行政当国。以朝诸侯。\\
\hspace*{2em}帝太甲居桐宫三年。悔过自责。反善。於是伊尹乃迎帝太甲而授之政。帝太甲修德。诸侯咸归殷。百姓以宁。伊尹嘉之。乃作太甲训三篇。褒帝太甲。称太宗。\\
\hspace*{2em}太宗崩。子沃丁立。帝沃丁之时。伊尹卒。既葬伊尹於亳。咎单遂训伊尹事。作沃丁。\\
\hspace*{2em}沃丁崩。弟太庚立。是为帝太庚。帝太庚崩。子帝小甲立。帝小甲崩。弟雍己立。是为帝雍己。殷道衰。诸侯或不至。\\
\hspace*{2em}帝雍己崩。弟太戊立。是为帝太戊。帝太戊立伊陟为相。亳有祥桑谷共生於朝。一暮大拱。帝太戊惧。问伊陟。伊陟曰。﹁臣闻妖不胜德。帝之政其有阙与。帝其修德。﹂太戊从之。而祥桑枯死而去。伊陟赞言于巫咸。巫咸治王家有成。作咸艾。作太戊。帝太戊赞伊陟于庙。言弗臣。伊陟让。作原命。殷复兴。诸侯归之。故称中宗。\\
\hspace*{2em}中宗崩。子帝中丁立。帝中丁迁于隞。河亶甲居相。祖乙迁于邢。帝中丁崩。弟外壬立。是为帝外壬。仲丁书阙不具。帝外壬崩。弟河亶甲立。是为帝河亶甲。河亶甲时。殷复衰。\\
\hspace*{2em}河亶甲崩。子帝祖乙立。帝祖乙立。殷复兴。巫咸任职。\\
\hspace*{2em}祖乙崩。子帝祖辛立。帝祖辛崩。弟沃甲立。是为帝沃甲。帝沃甲崩。立沃甲兄祖辛之子祖丁。是为帝祖丁。帝祖丁崩。立弟沃甲之子南庚。是为帝南庚。帝南庚崩。立帝祖丁之子阳甲。是为帝阳甲。帝阳甲之时。殷衰。\\
\hspace*{2em}自中丁以来。废適而更立诸弟子。弟子或争相代立。比九世乱。於是诸侯莫朝。\\
\hspace*{2em}帝阳甲崩。弟盘庚立。是为帝盘庚。帝盘庚之时。殷己都河北。盘庚渡河南。复居成汤之故居。乃五迁。无定处。殷民咨胥皆怨。不欲徙。盘庚乃告谕诸侯大臣曰。﹁昔高后成汤与尔之先祖俱定天下。法则可修。舍而弗勉。何以成德。﹂乃遂涉河南。治亳。行汤之政。然後百姓由宁。殷道复兴。诸侯来朝。以其遵成汤之德也。\\
\hspace*{2em}帝盘庚崩。弟小辛立。是为帝小辛。帝小辛立。殷复衰。百姓思盘庚。乃作盘庚三篇。帝小辛崩。弟小乙立。是为帝小乙。\\
\hspace*{2em}帝小乙崩。子帝武丁立。帝武丁即位。思复兴殷。而未得其佐。三年不言。政事决定於冢宰。以观国风。武丁夜梦得圣人。名曰说。以梦所见视群臣百吏。皆非也。於是乃使百工营求之野。得说於傅险中。是时说为胥靡。筑於傅险。见於武丁。武丁曰是也。得而与之语。果圣人。举以为相。殷国大治。故遂以傅险姓之。号曰傅说。\\
\hspace*{2em}帝武丁祭成汤。明日。有飞雉登鼎耳而呴。武丁惧。祖己曰。﹁王勿忧。先修政事。﹂祖己乃训王曰。﹁唯天监下典厥义。降年有永有不永。非天夭民。中绝其命。民有不若德。不听罪。天既附命正厥德。乃曰其奈何。呜呼。王嗣敬民。罔非天继。常祀毋礼于弃道。﹂武丁修政行德。天下咸驩。殷道复兴。\\
\hspace*{2em}帝武丁崩。子帝祖庚立。祖己嘉武丁之以祥雉为德。立其庙为高宗。遂作高宗肜日及训。\\
\hspace*{2em}帝祖庚崩。弟祖甲立。是为帝甲。帝甲淫乱。殷复衰。\\
\hspace*{2em}帝甲崩。子帝廪辛立。帝廪辛崩。弟庚丁立。是为帝庚丁。帝庚丁崩。子帝武乙立。殷复去亳。徙河北。\\
\hspace*{2em}帝武乙无道。为偶人。谓之天神。与之博。令人为行。天神不胜。乃僇辱之。为革囊。盛血。昂而射之。命曰﹁射天﹂。武乙猎於河渭之间。暴雷。武乙震死。子帝太丁立。帝太丁崩。子帝乙立。帝乙立。殷益衰。\\
\hspace*{2em}帝乙长子曰微子启。启母贱。不得嗣。少子辛。辛母正后。辛为嗣。帝乙崩。子辛立。是为帝辛。天下谓之纣。\\
\hspace*{2em}帝纣资辨捷疾。闻见甚敏。材力过人。手格猛兽。知足以距谏。言足以饰非。矜人臣以能。高天下以声。以为皆出己之下。好酒淫乐。嬖於妇人。爱妲己。妲己之言是从。於是使师涓作新淫声。北里之舞。靡靡之乐。厚赋税以实鹿台之钱。而盈钜桥之粟。益收狗马奇物。充仞宫室。益广沙丘苑台。多取野兽蜚鸟置其中。慢於鬼神。大聚乐戏於沙丘。以酒为池。县肉为林。使男女裸相逐其间。为长夜之饮。\\
\hspace*{2em}百姓怨望而诸侯有畔者。於是纣乃重刑辟。有炮烙之法。以西伯昌。九侯。鄂侯为三公。九侯有好女。入之纣。九侯女不喜淫。纣怒。杀之。而醢九侯。鄂侯争之彊。辨之疾。并脯鄂侯。西伯昌闻之。窃叹。崇侯虎知之。以告纣。纣囚西伯羑里。西伯之臣闳夭之徒。求美女奇物善马以献纣。纣乃赦西伯。西伯出而献洛西之地。以请除炮烙之刑。纣乃许之。赐弓矢斧钺。使得征伐。为西伯。而用费中为政。费中善谀。好利。殷人弗亲。纣又用恶来。恶来善毁谗。诸侯以此益疏。\\
\hspace*{2em}西伯归。乃阴修德行善。诸侯多叛纣而往归西伯。西伯滋大。纣由是稍失权重。王子比干谏。弗听。商容贤者。百姓爱之。纣废之。及西伯伐饥国。灭之。纣之臣祖伊闻之而咎周。恐。奔告纣曰。﹁天既讫我殷命。假人元龟。无敢知吉。非先王不相我後人。维王淫虐用自绝。故天弃我。不有安食。不虞知天性。不迪率典。今我民罔不欲丧。曰﹃天曷不降威。大命胡不至﹄。今王其柰何。﹂纣曰。﹁我生不有命在天乎。﹂祖伊反。曰。﹁纣不可谏矣。﹂西伯既卒。周武王之东伐。至盟津。诸侯叛殷会周者八百。诸侯皆曰。﹁纣可伐矣。﹂武王曰。﹁尔未知天命。﹂乃复归。\\
\hspace*{2em}纣愈淫乱不止。微子数谏不听。乃与大师。少师谋。遂去。比干曰。﹁为人臣者。不得不以死争。﹂乃彊谏纣。纣怒曰。﹁吾闻圣人心有七窍。﹂剖比干。观其心。箕子惧。乃详狂为奴。纣又囚之。殷之大师。少师乃持其祭乐器奔周。周武王於是遂率诸侯伐纣。纣亦发兵距之牧野。甲子日。纣兵败。纣走入。登鹿台。衣其宝玉衣。赴火而死。周武王遂斩纣头。县之大白旗。杀妲己。释箕子之囚。封比干之墓。表商容之闾。封纣子武庚。禄父。以续殷祀。令修行盘庚之政。殷民大说。於是周武王为天子。其後世贬帝号。号为王。而封殷后为诸侯。属周。\\
\hspace*{2em}周武王崩。武庚与管叔。蔡叔作乱。成王命周公诛之。而立微子於宋。以续殷后焉。\\
\hspace*{2em}太史公曰。余以颂次契之事。自成汤以来。采於书诗。契为子姓。其後分封。以国为姓。有殷氏。来氏。宋氏。空桐氏。稚氏。北殷氏。目夷氏。孔子曰。殷路车为善。而色尚白。\\
% ===== 索隐述赞（司马贞《史记索隐》）· 统一样式 =====
% 题头「索隐述赞曰」五字并入反白框；正文每四句一列，定宽盒 \hfil 居列中；
% \fontsize{24}{41} 列距对齐全书界行（121+41k）；\clearpage 使述赞独占一页。
\clearpage
\shuzanhead{索隐述赞曰}\linebreak
{\shuzanbody\fontsize{24}{41}\selectfont
\hbox to \hsize{\hfil 简狄吞乙。是为殷祖。玄王启商。伊尹负俎。\hfil}\linebreak
\hbox to \hsize{\hfil 上开三面。下献九主。旋师秦卷。继相臣扈。\hfil}\linebreak
\hbox to \hsize{\hfil 迁嚣圯耿。不常厥土。武乙无道。祸因射天。\hfil}\linebreak
\hbox to \hsize{\hfil 帝辛淫乱。拒谏贼贤。九侯见醢。炮烙兴焉。\hfil}\linebreak
\hbox to \hsize{\hfil 黄钺斯仗。白旗是悬。哀哉琼室。殷祀用迁。\hfil}%
\par}
\end{document}
